\documentclass[a4paper,11pt]{article}
\usepackage[utf8]{inputenc}
\usepackage[T1]{fontenc}
\usepackage[french]{babel}
\usepackage[left=2.5cm,right=2cm,top=2cm,bottom=1cm]{geometry}
\usepackage{hyperref}
\usepackage{xcolor}
\usepackage{fontawesome5}
\usepackage{lmodern}

% --- Couleurs ---
\definecolor{primary}{RGB}{0, 70, 135}

% --- Configuration des liens ---
\hypersetup{
    colorlinks=true,
    linkcolor=primary,
    urlcolor=primary,
}

\begin{document}

% --- EN-TÊTE EXPÉDITEUR ---
\noindent
\textbf{Loïc Aron MBASSI EWOLO} \\
Élève-Ingénieur ENSEA / Polytechnique Yaoundé \\
Cergy, France \\
\faPhone\ (+33) 7 59 52 33 98 \\
\faEnvelope\ \href{mailto:loicaron.mbassiewolo@ensea.fr}{loicaron.mbassiewolo@ensea.fr} \\
\faLinkedin\ \href{https://www.linkedin.com/in/loic-aron-mbassi-ewolo-3942a9246/}{LinkedIn} \quad
\faGlobe\ \href{https://mbassiloic.tech}{mbassiloic.tech}

\vspace{0.6cm}

\noindent

% --- DESTINATAIRE ---
\noindent
\textbf{CFL - Société Nationale des Chemins de Fer Luxembourgeois} \\
Service Recrutement et Accompagnement Carrière \\
À l'attention de Mme Kristina KAUT \\
Luxembourg

\vspace{0.4cm}

\hfill Cergy, le 23 janvier 2026

\vspace{0.4cm}

\noindent
\textbf{Objet : Candidature au stage Ingénieur Software IoT - Service Trains et Matériel}

\vspace{0.3cm}

\noindent
Madame,

\vspace{0.2cm}

Actuellement élève-ingénieur en deuxième année à l'ENSEA (Cergy) dans le cadre d'un double diplôme avec l'École Polytechnique de Yaoundé, je suis à la recherche d'un stage de 3 à 5 mois. Votre offre au sein du groupe CBM (Condition-Based Maintenance) a retenu mon attention car elle combine plusieurs domaines qui me motivent : l'IoT, l'analyse de données et le Machine Learning appliqués à un contexte industriel concret.

\vspace{0.3cm}

\textbf{Pourquoi cette offre m'intéresse :}

\vspace{0.2cm}

La mission de créer des tableaux de bord interactifs pour traiter les signaux d'alarme des trains me parle directement. Lors d'un hackathon (IndabaX Africa), j'ai développé un dashboard avec Streamlit et Flask pour visualiser des données médicales en temps réel. Cette expérience m'a appris à structurer l'information pour qu'elle soit exploitable rapidement par les utilisateurs. J'ai également travaillé sur des tableaux de bord pour une plateforme de télémédecine pendant deux ans, ce qui m'a familiarisé avec les problématiques de données en temps réel et d'ergonomie.

\vspace{0.3cm}

\textbf{Ce que je peux apporter :}

\vspace{0.2cm}

Sur le volet technique, je maîtrise Python et JavaScript, et j'ai une expérience pratique avec XML dans le cadre d'un projet de génération de code à partir de diagrammes UML. Concernant le Machine Learning, j'ai participé à un projet de maintenance prédictive (Urban Waste Management) où nous avons analysé des données de capteurs pour établir des corrélations entre historiques de pannes et mesures terrain. Ce projet, qui a remporté le premier prix au hackathon CITS, m'a sensibilisé aux enjeux de la maintenance conditionnelle.

\vspace{0.2cm}

J'ai également effectué un stage au laboratoire IoT de mon école, où j'ai travaillé sur l'optimisation de modèles ML pour du hardware embarqué (TinyML). Cette expérience m'a permis de comprendre les contraintes liées aux systèmes temps réel et à la collecte de données capteurs.

\vspace{0.3cm}

\textbf{Mes limites et ma motivation à apprendre :}

\vspace{0.2cm}

Je n'ai pas d'expérience directe dans le domaine ferroviaire ni avec les systèmes de GMAO. Cependant, je suis curieux d'apprendre et je m'adapte rapidement à de nouveaux environnements techniques. Ma formation actuelle en traitement du signal et systèmes embarqués à l'ENSEA me donne des bases solides pour comprendre les problématiques de signaux d'alarme et de capteurs.

\vspace{0.3cm}

Je suis disponible dès avril 2026 et mobile pour m'installer au Luxembourg pendant la durée du stage. Je serais heureux de pouvoir échanger avec vous sur cette opportunité et de vous présenter plus en détail mes réalisations.

\vspace{0.5cm}

\noindent
Je vous prie d'agréer, Madame, l'expression de mes salutations distinguées.

\vspace{0.5cm}

\noindent
\textbf{Loïc Aron MBASSI EWOLO}

\end{document}
